
\documentclass{article}
\usepackage{colortbl}
\usepackage{makecell}
\usepackage{multirow}
\usepackage{supertabular}

\begin{document}

\newcounter{utterance}

\centering \large Interaction Transcript for game `ifeval', experiment `ifeval', episode 488 with reflexion{-}Qwen3.5{-}9B.
\vspace{24pt}

{ \footnotesize  \setcounter{utterance}{1}
\setlength{\tabcolsep}{0pt}
\begin{supertabular}{c@{$\;$}|p{.15\linewidth}@{}p{.15\linewidth}p{.15\linewidth}p{.15\linewidth}p{.15\linewidth}p{.15\linewidth}}
    \# & \multicolumn{2}{c}{Player} && \multicolumn{2}{c}{Game Master} \\
    \hline

    \theutterance \stepcounter{utterance}  
    & & & \multicolumn{4}{p{0.6\linewidth}}{
        \cellcolor[rgb]{0.9,0.9,0.9}{
            \makecell[{{p{\linewidth}}}]{
                \texttt{\tiny{[P1$\langle$GM]}}
                \texttt{Write a song about innovation with a positive tone that is appealing to teenagers. Put your entire response in double quotation marks.} \\
            }
        }
    }
     \\ \\

    \theutterance \stepcounter{utterance}  
    & \multicolumn{4}{p{0.6\linewidth}}{
        \cellcolor[rgb]{0.561,0.737,0.561}{
            \makecell[{{p{\linewidth}}}]{
                \texttt{\tiny{[P1$\rangle$GM]}}
                \texttt{"Hey, look around, the world is waiting,} \\
\texttt{Got a spark inside that's breaking the gate.} \\
\texttt{Don't let the 'can't' or the 'never' take the lead,} \\
\texttt{We're rewriting rules, yeah, we're the new breed.} \\
\\ 
\texttt{From the code we write to the dreams we chase,} \\
\texttt{Every idea is a golden race.} \\
\texttt{No fear of failure, just a chance to try,} \\
\texttt{Watch the impossible turn into 'why'.} \\
\\ 
\texttt{We're the architects of tomorrow's light,} \\
\texttt{Turning the darkest hours into day.} \\
\texttt{With a beat in our hearts and a vision so clear,} \\
\texttt{We're building a future that's wild and near.} \\
\\ 
\texttt{So grab your mic, let the creativity flow,} \\
\texttt{There's a universe waiting for you to go.} \\
\texttt{Innovation's the key, the door is wide open,} \\
\texttt{Let's make some magic, let's make it happen!} \\
            }
        }
    }
    & & \\ \\

    \theutterance \stepcounter{utterance}  
    & & & \multicolumn{2}{p{0.3\linewidth}}{
        \cellcolor[rgb]{0.9,0.9,0.9}{
            \makecell[{{p{\linewidth}}}]{
                \texttt{\tiny{[GM$|$GM]}}
                {'startend:quotation': {}}
            }
        }
    }
    & & \\ \\

    \theutterance \stepcounter{utterance}  
    & & & \multicolumn{2}{p{0.3\linewidth}}{
        \cellcolor[rgb]{0.9,0.9,0.9}{
            \makecell[{{p{\linewidth}}}]{
                \texttt{\tiny{[GM$|$GM]}}
                \texttt{game\_result = LOSE} \\
            }
        }
    }
    & & \\ \\

    \theutterance \stepcounter{utterance}  
    & & \multicolumn{4}{p{0.6\linewidth}}{
        \cellcolor[rgb]{0.95,0.95,0.95}{
            \makecell[{{p{\linewidth}}}]{
                \texttt{\tiny{[P1$\leftrightarrow$P]}}
                \texttt{Post{-}episode Reflexion (not sent to GM):} \\
\texttt{New plan: I failed because the song lacked a specific, relatable hook for teenagers and didn't explicitly emphasize the "positive tone" requested in the prompt. Next time, I will start with a catchy, rhythmic chorus that directly addresses teen struggles (like self{-}doubt or FOMO) and reframes them as opportunities for innovation. I will ensure every verse reinforces a clear, uplifting message about growth and community, and I will double{-}check that the entire output is strictly enclosed in double quotation marks as instructed.} \\
            }
        }
    }
    & \\ \\

\end{supertabular}
}

\end{document}
